% Figure: schematic of a two-mass, two-spring chain and its two modes.
% Two masses on a line; arrows show in-phase (mode 1) and
% out-of-phase (mode 2) motion. Springs drawn as zigzag polylines so
% the figure needs no decoration libraries beyond the preamble defaults.
\begin{figure}[htbp]
\centering
\begin{tikzpicture}[
  mass/.style={draw, fill=engblue!12, minimum width=0.9cm, minimum height=0.9cm},
  arr/.style={-{Stealth}, very thick},
]
% spring as a zigzag from (xa,y) to (xb,y); fixed 6-segment shape.
\newcommand{\spring}[3]{%
  \draw[enggray, thick] (#1,#3) -- (#1+0.12,#3)
    -- (#1+0.22,#3+0.13) -- (#1+0.40,#3-0.13)
    -- (#1+0.58,#3+0.13) -- (#1+0.76,#3-0.13)
    -- (#1+0.94,#3+0.13) -- (#1+1.08,#3) -- (#2,#3);}
% === top row: mode 1, in phase ===
\node[font=\small\bfseries, anchor=west] at (-1.0,2.4) {Mode 1: in phase (low frequency)};
\draw[enggray, thick] (-0.7,1.45) -- (-0.7,2.05);
\spring{-0.7}{0.6}{1.75}
\node[mass] (m1a) at (1.05,1.75) {$m_1$};
\spring{1.5}{2.7}{1.75}
\node[mass] (m2a) at (3.15,1.75) {$m_2$};
\draw[arr, engred] (1.05,1.0) -- (1.55,1.0);
\draw[arr, engred] (3.15,1.0) -- (3.65,1.0);
% === bottom row: mode 2, out of phase ===
\node[font=\small\bfseries, anchor=west] at (-1.0,0.1) {Mode 2: out of phase (high frequency)};
\draw[enggray, thick] (-0.7,-0.85) -- (-0.7,-0.25);
\spring{-0.7}{0.6}{-0.55}
\node[mass] (m1b) at (1.05,-0.55) {$m_1$};
\spring{1.5}{2.7}{-0.55}
\node[mass] (m2b) at (3.15,-0.55) {$m_2$};
\draw[arr, engred] (1.05,-1.3) -- (1.55,-1.3);
\draw[arr, engblue] (3.15,-1.3) -- (2.65,-1.3);
\end{tikzpicture}
\caption[The two modes of a two-mass chain]{The two natural modes of
a two-mass, two-spring chain anchored to a wall. In the first
(low-frequency) mode the masses move in phase, in the same direction
at the same instant; the connecting spring between them barely
stretches, so the restoring force is small and the frequency is low.
In the second (high-frequency) mode the masses move out of phase, in
opposite directions; the connecting spring is maximally strained, the
restoring force is large, and the frequency is high. Every linear
multi-degree-of-freedom system decomposes into such modes, each an
independent single-degree-of-freedom oscillator in its own modal
coordinate. The mode shapes, not the physical coordinates, are the
natural variables of the dynamics.}
\label{fig:vol03ch06:2dof-modes}
\end{figure}
